\section{位移的物理账本（一）：地面交通的代价}

\subsection{动力平台、位移系统与赛道误判}

机制可以在单个环节里被辨认，却只有在完整系统中才能被定价。资源、时间与拓扑三类约束，在单项技术里往往表现为孤立参数；一旦进入真实运行的系统，它们立刻互相置换，把代价从一处搬到另一处。局部指标的改善与系统总负担的下降，是两件不同的事，而后者只能在整个系统上称量。要看清这个区别，需要一个足够沉重、足够长期、也足够耦合的对象。

地面交通系统是技术黏滞的绝佳标本。它同时包含三类约束：道路、桥梁、停车场、车辆、电网、充电网络、维修体系、保险制度、交通法规、驾驶习惯、城市空间和家庭生活，全都被锁在同一个系统里。交通从来不是单车技术问题，而是一个高度物质化、长期沉积、高度耦合的状态转移问题。换一辆车容易，改变位移系统很难；替换动力平台容易，重构城市空间很难。

\subhead{动力平台与位移系统}

过去几年，“新能源电车”常被包装成未来交通的核心答案。这套叙事的问题不在于电动车没有技术价值，而在于它把局部电动化、智能座舱、产业政策、资本故事、消费升级和民族工业情绪混成了一种神话。车企说智能，政策说绿色，资本说增长，媒体说弯道超车，消费者说科技感。可这些词叠在一起，并不会自动降低交通系统的总黏滞。

承认车辆端的技术优势，并不改变这笔账。电机效率高、响应快、结构相对简洁，电控系统也更容易接入数字平台。这些都是真实优势。但交通系统的瓶颈从来不只是电机效率。电机再轻巧，电池仍然笨重；电控再精细，车辆仍然占地；加速再漂亮，道路仍然拥堵；座舱再智能，大多数私家车仍然长期闲置。把电动车说成交通革命，本质上是把车辆内部的局部优化，夸大成整个交通系统的跃迁。

电动车首先是动力平台替换，不是交通系统重构。它把发动机、油箱、变速箱和加油体系，替换成电池、电机、电控、热管理和充电体系。内燃机的复杂性下降了，电池系统的复杂性上升了；尾气从车尾消失了，能源、矿产、制造、回收和电网压力转移到了系统后台。一个环节变轻，另一个环节可能变重；一个机械结构变简单，基础设施和资源系统可能变复杂。如果只看车辆端宣传图像，而不看电池、矿产、电网、道路、停车、保险、折旧和利用率，就仍然是单点技术幻觉。

转型路径要比较完整生命周期。沿用既有结构可以节省近期投资，也可能延长燃料、维护与排放支出；新建基础设施有成本，也可能降低长期总成本。力学上的最小作用量原理不替这项比较作决定。应当核算的是同等服务下的能源、材料、费用、排放、空间与风险。

充电、配电、热管理和电池回收构成新的支出，燃油开采、炼制、运输、加油与发动机维护也有原来的支出。IEA的2026年生命周期比较仍给出电动车的排放收益，同时列明制造与供电条件。\cite{iea2026evlifecycle} 电动化可以形成净改进，拥堵、停车与城市可达性则仍需另作改造。

这正是技术黏滞论要强调的判断：\key{降低某一局部黏滞，不等于降低系统总黏滞。}电动车可能降低了车辆运行中的部分能量转换损耗，也可能降低某些维护成本，但它没有自动解决交通系统中更深的时间刚性和拓扑耦合。道路仍然拥堵，车辆仍然闲置，停车仍然占地，城市仍然围绕汽车扩张，家庭仍然被迫购车，通勤仍然吞噬时间，事故责任仍然存在。车变了，位移系统未必变了。要把这个判断变成可用的工具，需要先确认被优化的对象究竟是什么：如果对象是车，动力平台替换已经足够；如果对象是位移，结论就会不同。

\subhead{总位移与有效位移}

交通的第一性对象不是车，而是位移。人、货物和服务必须在空间中重新分布。人要抵达工作、教育、医疗、社交和家庭场景；货物要穿过生产、仓储、干线运输、末端配送和消费节点；服务要在正确时间抵达正确地点。车只是位移载体之一。步行、自行车、公交、地铁、铁路、货车、船运、自动驾驶运营车和末端配送系统，都是不同尺度下的位移工具。交通系统真正要优化的，不是每辆车有多先进，而是必要位移如何以更低物理代价、更高可靠性和更强韧性完成。

承载效率不由车辆的动力形式决定，它受密度与速度的关系约束。一条道路断面在单位时间内通过的车辆数由下式给出：
\begin{equation}
q=\rho v
\label{eq:ch25-flow}
\end{equation}
其中 $q$ 是流量，量纲为辆每小时；$\rho$ 是车辆密度，量纲为辆每千米；$v$ 是该断面的平均速度，量纲为千米每小时。这个关系本身只是量的定义，它的实质内容在于 $v$ 随 $\rho$ 上升而下降：车流越密，车头间距越小，速度越低。二者相乘的结果是 $q$ 在某一临界密度处取到最大值，超过该密度后流量反而下降，路网进入拥堵分支\cite{helbing2001traffic}。动力平台替换可以改变单车的能耗、排放与噪声，却不会移动这条曲线，因为曲线的位置由车头间距、驾驶人反应时间和道路几何决定。

更准确地说，交通系统的第一性问题，是在多重约束下完成有效位移的最小作用量优化。人和货物必须从 A 点到 B 点，但真正需要最小化的不是单次行驶时间，也不是单车能耗，而是整个系统为完成这次位移所支付的总作用量。这里的“总作用量”不是狭义物理公式，而是对交通系统全部代价的概括：车辆制造与折旧、能源消耗、道路和停车占地、驾驶注意力、拥堵外部性、事故风险、维护成本、等待时间、换乘摩擦、排放、噪声和系统韧性损耗。所谓好交通，不是让某一种车辆在某一项指标上达到极值，而是在资源、时间、空间、风险和自由度之间找到更低总代价的位移路径。

因此，需要区分“总位移”和“有效位移”。并不是所有移动都代表进步。过长通勤、职住分离、空驶车辆、重复配送、低效换乘、被迫跨区就医、长期寻找停车位，都是系统结构低效造成的位移。一个交通系统如果只是让这些无效移动变得更电动、更安静、更智能，它并没有真正变好。更好的交通，不只是提高移动速度，而是减少无效移动，提高必要位移的承载效率。

从这个角度看，私家车模式是一个典型的局部最优。私家车当然有价值。它提供自由、弹性、隐私和可达性。对低密度地区、公共交通不足地区、家庭接送、老人小孩照护和复杂日程来说，私家车不是简单奢侈品，而是现实生活的基础接口。很多人依赖汽车，不是因为他们被广告洗脑，而是因为居住、工作、商业、学校、医院和服务网络已经围绕汽车组织起来。没有车，生活就会变得困难。

但私家车也是高黏滞结构。它利用率极低，大部分时间停放；它占用道路和停车空间；它要求个人承担驾驶、维护、保险、补能和事故风险；它把人的注意力锁进驾驶任务；它推动城市空间向低密度、长距离、车依赖方向扩展。更重要的是，一旦社会围绕私家车建成，离开它就很难。道路、郊区、商场、停车位、贷款、保险、家庭时间表和生活习惯共同形成结构记忆。私家车模式不是最优系统，但它是一个稳定的局部最优。局部最优之所以稳定，是因为围绕它形成的依赖关系彼此加固。依赖关系不改变，替换动力平台就不会改变系统结构。

\subhead{赛道误判}

这也是为什么单纯“油车换电车”改变不了交通底层结构。它保留了私家车模式的大部分拓扑：个人占有、低利用率、停车占地、道路拥堵、家庭补能、分散维护、事故责任和城市扩张。区别只是动力平台换了。系统仍然让一两个人坐在一两吨机器里，占用道路、停车和能源网络，完成许多本可以用更轻、更集约、更少位移的方式完成的任务。电动车解决了部分动力问题，却没有解决“为什么要移动这么多车”的问题。

这里需要把两组关系分清楚。电动车和燃油车，是动力平台之争；私家车和自动驾驶运营车，才是位移组织方式之争。电动化当然重要，它很可能是自动驾驶运营车的基础条件，因为运营车队需要更低维护成本、更容易控制的动力系统和更适合集中补能的能源接口。真正的问题不是电动车对不对，而是电动化被锁进了旧的私家车范式。把电动私家车保有量本身当作未来交通目标，就是把工具误认为终点。

更糟的是，许多新能源车叙事并不反思这个问题，反而把它包装得更高级。更大的屏幕、更强的算力、更重的电池、更复杂的座舱、更夸张的加速、更密集的传感器，被讲成“智能化”和“进化”。但如果这些东西只是让一辆低利用率私家车变得更重、更贵、更难维护、更依赖后台系统，它们未必是在降低黏滞，也可能是在用科技感遮蔽新的黏滞。车越来越像移动电子消费品，交通系统却仍然像旧时代一样沉重。

“弯道超车”这个叙事最危险的地方，是它可能误判了赛道。它默认未来交通仍然以大规模私人购车为中心，只是把油车换成电车，把发动机换成电池和电机，把机械座舱换成智能座舱。于是政策、补贴、路权、牌照、产业投资和舆论资源都围绕新能源私家车展开，仿佛只要在这一代私人车辆平台上取得优势，就等于占据了未来交通的制高点。

可是，更接近未来交通主形态的，可能不是每个家庭拥有一辆更先进的电动车，而是自动驾驶运营车大规模覆盖日常出行：无人出租车、无人通勤车、自动驾驶接驳车、自动配送车，以及在固定区域和固定线路内运行的运营车辆。它们一旦成熟，会直接吃掉大量今天由私家车承担的场景：上下班通勤、老人小孩接送、短途办事、城市内移动、机场车站接驳、夜间出行和末端配送。到那时，竞争的核心不再是谁卖出更多私人电动车，而是谁能用更少车辆、更高利用率、更低等待时间、更稳定安全边界，完成更多日常位移。从这个角度看，强行推动新能源私家车的政策叙事有一种反讽：它以未来的名义强化旧范式。它打击燃油车，却未必打击私家车依赖；它鼓励电动化，却未必鼓励位移效率；它扶持车辆制造，却未必重构城市交通；它制造“先进产业”的兴奋感，却可能继续加固低利用率、高占地、高折旧、高资源消耗的私人车辆结构。

如果未来交通的核心转向自动驾驶运营车，那么今天许多被称为“弯道超车”的政策，可能只是把旧赛道铺得更宽，最后却发现真正的赛道已经转向了运营化、高利用率车辆网络和有效位移成本竞争。所谓把车道超没了，正是这个意思。一个产业以为自己在弯道超车，结果未来的竞争对象不再是另一辆私人车，而是自动驾驶运营车对私家车场景的系统性吞并。私家车仍会存在，但它不再天然占据交通中心；车企仍会重要，但单车销售不再必然是最大价值入口；电池、电机、车身和座舱仍有技术含量，但决定交通未来的可能是运营能力、车辆利用率、安全边界、维护体系、责任链条和单位有效位移成本。真正的超车，不是把旧车道上的私人车造得更快，而是让大量日常位移不再需要那条旧车道。如果运营车是竞争的下一个入口，那么它的技术形态就决定了黏滞会被搬到哪里。

运营化之所以可能改变总黏滞，是因为它同时作用于三项来源，而三项来源的响应方向并不相同。车辆以运营方式集中使用，压低的是资源势垒：同样数量的位移由更少的车辆完成，制造、折旧、停车占地和补能设施的绝对需求随之下降。车辆可以集中维护、连续运行、动态调度，压低的是拓扑耦合：位移不再必须由每个家庭自备一台机器去对接道路网络，车辆与空间的绑定关系松动。但运营化同时抬高时间刚性：一支在开放道路上连续运行的无人车队，需要更长的验证周期、更完整的责任链条和更保守的安全边界。三项来源不会同时下降，只会重新分布。判断自动驾驶的技术形态，不能只问它是否更智能，还要问它把黏滞搬到了哪一维。

\subsection{运营化、游艇化与多层系统的韧性}

\subhead{复杂度棘轮}

自动驾驶不能被神话化。它不是智能直接降临交通系统，而是典型的覆盖层策略：在旧道路、旧车辆、旧法规、旧保险、旧行为习惯和旧城市空间上，叠加传感器、地图、模型、算力、远程接管、日志、审计和责任机制。它试图让一个为人类身体和人类判断建成的世界，变得可机器读取、可机器预测、可机器执行。它不是消除黏滞，而是把黏滞转移到车辆复杂度、场景标定、长尾验证、基础设施接口和责任链条上。

这就是复杂度棘轮。为了让机器在旧世界中行动，系统必须不断增加新层：摄像头、雷达、激光雷达、高精地图、车路协同、云端训练、仿真平台、远程运营、数据闭环、安全员、法规试点、保险安排、事故调查。每一层都降低了某些操作摩擦，却增加了系统耦合和维护负担。单车看起来更聪明，交通系统却不一定更简单。智能覆盖层越厚，后台越重。

自动驾驶最难的场景，是人机混行的开放道路。纯人类驾驶是一种旧秩序。它不完美，但人类司机、行人、骑行者、交警、保险、道路习惯和责任规则已经共同形成某种可运行的生态。高度标准化的纯机器系统也可能稳定，例如封闭园区、港口、矿区、固定路线和专用道路。最困难的是过渡态：人类司机、机器车辆、行人、非标准道路、临时施工、天气变化、地方习惯和责任争议混在一起。这里不是单车驾驶问题，而是一个高耦合混合系统问题。

L5 自动驾驶的困难也不能只归结为模型不够强。开放道路包含大量残差：模糊手势、临时障碍、非规则行为、路面磨损、施工改道、动物、小孩、老人、骑行者、事故现场、地方性默契和不可预测的社会互动。人类身体和经验可以在很多非标准情境中低成本处理残差，机器则必须把这些残差显式工程化：传感、识别、预测、决策、验证、记录、追责。这正是身体封装与复杂度显影的差异：人把大量现实细节压在身体、经验和默契之中，机器则必须把这些隐含处理拆成可感知、可计算、可验证、可追责的工程结构。

所以，自动驾驶真正容易落地的地方，往往不是最开放、最混乱、最“像人”的道路，而是运行域更清晰的场景：固定区域、固定线路、封闭园区、港口、矿区、仓储、机场接驳、低速配送、专用车道、V2X 协同环境。它们不是让机器完全像人一样适应旧世界，而是把世界局部低熵化，让机器不必重建完整的人类身体和社会经验。很多自动化真正落地，不是因为机器突然拥有通用智能，而是因为环境被改造成了机器可承受的形态。落地条件的分析同时给出另一条结构线索：可以被改造的既可以是环境，也可以是车辆的组织方式，而后者的影响范围更大。

\subhead{运营化与游艇化}

运营车队可能通过提高利用率、集中维护和调度减少部分私有资产需求，也会产生空驶、清洁、支持、平台与集中失效风险。是否降低总黏滞，要比较完整服务：等待时间、覆盖、费用、车辆公里、道路占用与可靠性。无人驾驶扩大了运营选择，并不自动把每一项系统指标推向更优。

这会改变私家车的社会角色。私家车不会消失，但它会从大众基础设施退回到消费品、兴趣品、身份品和特殊场景资产。今天多数人拥有汽车，是因为生活结构迫使他们拥有汽车；未来如果自动驾驶运营车覆盖多数日常场景，日常位移不再依赖私人占有车辆，那么私家车就会逐渐游艇化。它仍然有人喜欢，仍然有人收藏，仍然有人为自由感、速度感、审美和身份支付高价，但它不再是多数人完成日常生活的必要基础设施。

“游艇化”是一种有条件的结构判断：如果运营服务足以覆盖多数日常需求，私人占有车辆可能更多承担兴趣、审美与特殊用途。覆盖、价格、隐私、残障服务、儿童出行、低密度地区与应急韧性，决定条件能否成立。要观察的是家庭是否实际减少车辆，以及总行驶量与服务质量怎样变化。

\subhead{多层系统与韧性}

未来交通必须是多层系统，而不是单一工具统治。短距离适合步行、自行车和轻型交通；高密度走廊适合公交、地铁和轨道；低密度弹性出行适合自动驾驶运营车；大宗货运适合铁路、船运和干线物流；末端配送需要小型化、自动化和人类残差处理。真正高效的交通，不是一种工具覆盖所有场景，而是不同位移需求匹配不同物理载体。

交通还必须保留韧性。极致效率不等于好系统。天气、事故、能源波动、设备故障、人口变化、突发事件和技术迭代，都会扰动位移网络。一个交通系统如果为了效率消灭所有冗余，正常状态下可能漂亮，一旦扰动出现就会迅速失稳。韧性来自多样性、替代路径和局部冗余。轨道交通高效但不够灵活，汽车灵活但利用率低，自行车轻便但受距离和天气限制，自动驾驶运营车可持续运行但依赖传感器、算法和基础设施。它们各有边界，也正因为彼此不同，才可能互相补足。

这就是交通系统中的黏滞管理。不是消灭移动，而是减少无效移动；不是追求单车最优，而是降低单位有效位移的总代价；不是用电动车替代油车就宣布未来到来，而是重新组织车辆、道路、能源、空间、运营、责任和韧性；不是把所有交通工具合并成一种终极机器，而是让不同工具在不同尺度上承担合适任务。交通的未来，不是车辆形态的单次升级，而是大量日常位移从私人车辆占有，转向自动驾驶运营车承接。电动化很重要，但它只是底座；自动驾驶很重要，但它的真正形态不应是更昂贵的无人私家车。真正的结构变化，是私家车基础设施地位的下降。

到那时，车仍然会存在。但许多私家车会像游艇一样存在：不再是日常生活的基础设施，而是少数人愿意为占有、审美、自由感和身份支付的高黏滞消费品。私家车地位的下降，仍然是一次发生在同一张网络上的重新分配。无论车辆归谁所有、由谁运营，位移都要经过道路、桥梁、隧道、路口和产权边界构成的二维拓扑。当某一段位移的地面路径本身过重时，降低车辆成本并不能缩短它。位移系统内部还能挤出多少效率，与能不能绕开这张网络，是两个不同的问题。
